\maketitle

% sections/00_frontmatter.tex
% Abstract + keywords + JEL. Title/authors/date assumed in main.tex (\maketitle).
% Bibitem keys used below: Bruegel2023, auclert2023energy. Align to bibliography.tex if named differently.

\begin{abstract}
\noindent European governments spent more than \euro{}600 billion shielding households from
the 2021--23 energy shock, mostly by suppressing retail energy prices
\citep{Bruegel2023}. We show that holding down the price also holds back the crisis-time adoption of
green durables, a margin the macroeconomic literature has left out. We add a discrete brown/green durable choice to the small open economy
HANK model of \citet{auclert2023energy}. An energy shock raises the relative price of brown
energy and triggers an adoption wave: the green share rises by about eight
percentage points at the peak, from a five-percent base. The switch is a lumpy
purchase financed against a borrowing constraint, so adoption concentrates among
unconstrained households. A retail price
cap freezes the brown price and eliminates the wave; an income transfer of equal fiscal
cost leaves the price free and preserves it, at comparable output stabilisation.
The same suppression shuts down the expenditure-switching gain a small open
economy draws from the shock. Two further results follow. An untargeted flat
transfer dominates one indexed to pre-crisis energy use, because homothetic energy
demand makes energy-indexed targeting regressive. A permanent carbon tax that
greens the steady state provides no first-order crisis insurance. Monetary policy
moves the transition only weakly, through the financing cost of the switch.
Governments that shield households should support incomes and leave the retail
price free.
\end{abstract}

\medskip
\noindent\textbf{Keywords:} energy shocks; green transition; durable adoption;
HANK; price caps; transfers; carbon taxation.

\smallskip
\noindent\textbf{JEL codes:} E21, E32, F41, H23, Q43, Q58.

\thispagestyle{empty}
% H23 added (externalities; redistribution) given the policy section; drop if the
% coauthors prefer the tighter set E21 E32 F41 Q43 Q58.
\clearpage
